\chapter{Programming quantum computers with Qiskit}\label{ch:qiskit}

\section{General remarks}
It is instructive to compare present-day quantum computers with the first general-purpose electronic computer, ENIAC (1945).\footnote{Before electronic computers, calculations relied on
analog devices such as the slide rule, used from the 17th century up to the 1970s.} ENIAC worked in decimal arithmetic and needed 36 vacuum tubes to store a single decimal digit. Its memory
consisted of 20 accumulators, the equivalent of modern registers, each holding a signed 10-digit number. Some of them were combined for more complex operations: four accumulators were used
for multiplication and five for division. Since $10^{10}\approx2^{33}$, one accumulator corresponds roughly to a 32-bit binary register, and the whole machine stored only a few hundred bits.
Even so, the accumulators alone required of the order of $20\times10\times36\approx10^4$ vacuum tubes.

Current quantum processors have a comparable number of (noisy) qubits. In terms of programming, quantum computing is therefore still at the stage of ENIAC, or even before it: a computer is in
effect designed with a specific hardware wiring for each problem. Programming is done at the lowest level, gate by gate, even though some higher-level codes are already available.

Several software frameworks exist for programming quantum computers (Qiskit, Cirq, PennyLane, Q\#, CUDA-Q, \dots). We use Qiskit, an open-source Python library developed by IBM. It can run
circuits both on simulators and on real quantum hardware. The examples below have been tested with Qiskit 2.x and Qiskit Aer 0.18.

\section{Building a circuit}
A circuit is built from registers:
\begin{lstlisting}
from qiskit import QuantumRegister, ClassicalRegister, QuantumCircuit

qreg = QuantumRegister(2, 'q')    # one quantum register with 2 qubits q0, q1
creg = ClassicalRegister(2, 'c')  # one classical register with 2 bits
qc = QuantumCircuit(qreg, creg)   # a circuit acting on these registers

qc.h(qreg[0])                     # Hadamard gate on q0
qc.cx(qreg[0], qreg[1])           # controlled-X (CNOT): control q0, target q1
qc.draw('mpl')                    # draw the circuit with matplotlib
\end{lstlisting}
The method \texttt{cx} adds a controlled-$X$ gate, i.e.\ the CNOT gate, with \texttt{qreg[0]} as control and \texttt{qreg[1]} as target. This circuit prepares the Bell state
$\ket{\Phi^+}$ (\cref{ch:entanglement}). The option \texttt{'mpl'} of the \texttt{draw} method produces a figure using the Python plotting library matplotlib.

\paragraph{Exercise: matrices by hand.} Before using a simulator, it is useful to compute the matrices of small circuits by hand, for example with the symbolic library SymPy. Given the CNOT
matrix \texttt{Cx\_sp}, construct the matrices \texttt{H\_sp}, \texttt{I\_sp}, \texttt{IH\_sp} $=I\otimes H$, and $U=(I\otimes H)\,\CNOT\,(I\otimes H)$:
\begin{lstlisting}
import sympy as sp
from sympy.physics.quantum import TensorProduct

H_sp = sp.Matrix([[1, 1], [1, -1]]) / sp.sqrt(2)
I_sp = sp.eye(2)
Cx_sp = sp.Matrix([[1, 0, 0, 0], [0, 1, 0, 0],
                   [0, 0, 0, 1], [0, 0, 1, 0]])
IH_sp = TensorProduct(I_sp, H_sp)
U_sp = sp.simplify(IH_sp * Cx_sp * IH_sp)
\end{lstlisting}
The result is $U=\diag(1,1,1,-1)$, the controlled-$Z$ gate. Since $HXH=Z$, the Hadamard gates on the target convert the controlled-$X$ into a controlled-$Z$. By hand: $I\otimes H$ is
block diagonal with two blocks $H$, and $\CNOT=\diag(I,X)$, so $U=\diag(HIH,HXH)=\diag(I,Z)$.

\section{Simulators}
Aer is the Qiskit package for high-performance simulation of quantum circuits, including noisy ones. It provides several simulator backends:
\begin{lstlisting}
from qiskit_aer import Aer
simulator = Aer.get_backend('qasm_simulator')
\end{lstlisting}
\begin{itemize}
  \item \texttt{qasm\_simulator}: simulates measurements shot by shot, as a real device would. It is ideal by default, but noise models can be added.
  \item \texttt{statevector\_simulator}: returns the ideal final state vector.
  \item \texttt{unitary\_simulator}: returns the ideal unitary matrix of the whole circuit.
\end{itemize}
In recent versions all three are special cases of the general \texttt{AerSimulator}, whose simulation method can be \texttt{'statevector'} (on CPU or GPU), \texttt{'density\_matrix'} (also on
GPU), \texttt{'unitary'}, \texttt{'matrix\_product\_state'}, \texttt{'stabilizer'}, and others.

\begin{keypoint}[: the cost of classical simulation]
A state vector of $n$ qubits has $2^n$ complex amplitudes, i.e.\ $16\times2^n$ bytes in double precision: 16\,GB for $n=30$, 16\,PB for $n=50$. A density matrix needs the square of
that. Brute-force simulation is therefore limited to about 45--50 qubits on the largest supercomputers. Matrix-product-state and stabilizer simulators can handle many more qubits for weakly
entangled or Clifford circuits, in agreement with the remark at the end of \cref{sec:parallelism}.
\end{keypoint}

\section{Transpiling and running a circuit}
Before running, the circuit is \emph{transpiled}. Transpilation translates it into the gates actually available on the backend, maps the logical qubits onto physical qubits respecting the
connectivity (adding SWAPs if needed), and optimizes it. Unlike compilation, which translates a program to a lower level of abstraction, this is a conversion at the same level, from a circuit
to an equivalent circuit:
\begin{lstlisting}
from qiskit import transpile
from qiskit.visualization import plot_histogram

qc.barrier()            # visual separator; also blocks optimizations across it
qc.measure(qreg, creg)  # measure the qubits into the classical bits

qc_t = transpile(qc, simulator,
                 basis_gates=['id', 'rz', 'sx', 'x', 'cx'])
job = simulator.run(qc_t, shots=10000)
counts = job.result().get_counts(qc_t)
plot_histogram(counts)
\end{lstlisting}
A barrier is a visual marker in the circuit drawing. It also prevents the transpiler from moving or merging gates across it. The result \texttt{counts} is a Python dictionary that gives the number
of times each outcome was obtained. For the Bell state prepared above, it is approximately \texttt{\{'00': 5000, '11': 5000\}}. The statistical fluctuation of each count is $\sqrt{Np(1-p)}=50$.
The basis gates \texttt{rz}, \texttt{sx}$=\sqrt X$, \texttt{x} suffice for any single-qubit gate (Euler decomposition, \cref{thm:euler}), e.g.\ $H=\e^{\ii\pi/4}R_z(\pi/2)\sqrt XR_z(\pi/2)$.

\section{State vector and unitary matrix}
With the state-vector simulator, applied to a circuit without the final measurements, the state of the circuit can be retrieved and displayed as a \LaTeX{} formula:
\begin{lstlisting}
from qiskit.visualization import array_to_latex

simulator = Aer.get_backend('statevector_simulator')
psi = simulator.run(qc_t).result().get_statevector(qc_t)
array_to_latex(psi)
\end{lstlisting}
Similarly, the unitary simulator gives the matrix of the whole circuit, which can be checked numerically, for example its determinant:
\begin{lstlisting}
import numpy as np
from numpy import linalg as LA

simulator = Aer.get_backend('unitary_simulator')
U = simulator.run(qc_t).result().get_unitary(qc_t, decimals=3)
U = np.asarray(U)      # convert to a NumPy array
print(LA.det(U))       # |det U| = 1, since U is unitary
\end{lstlisting}

\section{Qubit ordering: big-endians and little-endians}\label{sec:endian}
Qiskit orders the qubits from right to left: in a basis state $\ket{q_{n-1}\dots q_1q_0}$, qubit $q_0$ is the least significant bit. This is the opposite of the convention used in these notes.
As a result, the matrices returned by Qiskit may appear with the roles of the qubits exchanged.

\paragraph{A little history.} In Jonathan Swift's \emph{Gulliver's Travels} (1726), the empire of Lilliput is torn by a bitter conflict about how a boiled egg should be opened. Traditionally, eggs
were broken at the larger end. However, the Emperor's grandfather cut his finger while doing so as a child, and his father then ordered all his subjects, under heavy penalties, to break their eggs
at the smaller end. The resistance of the Big-Endians led to several rebellions. The exiled rebels found refuge in the neighbouring empire of Blefuscu, and a long war followed. Swift's satire
targets the religious quarrels of his time. The point is that the dispute is passionate, but the choice itself is completely arbitrary.

In 1980, the computer scientist Danny Cohen borrowed Swift's terms in a famous note, \emph{On Holy Wars and a Plea for Peace} \cite{Cohen1981}. Cohen was writing about the order in which computers
store and transmit the bytes of a number. A 32-bit integer such as \texttt{0x12345678} occupies four bytes, and there are two natural ways to arrange them at increasing addresses:
\begin{center}
\begin{tabular}{@{}lcccc@{}}\toprule
Address & $n$ & $n+1$ & $n+2$ & $n+3$\\ \midrule
Big-endian & 12 & 34 & 56 & 78\\ Little-endian & 78 & 56 & 34 & 12\\ \bottomrule
\end{tabular}
\end{center}
A big-endian machine stores the most significant byte first, the way we write numbers on paper. A little-endian machine stores the least significant byte first. Little-endian ordering simplifies some
arithmetic operations, since addition proceeds from the least significant digit; big-endian ordering is easier for humans to read in a memory dump. As in Lilliput, neither is fundamentally better,
and both have been used: the Intel x86 family and, in practice, the ARM processors of today's smartphones and Apple silicon Macs are little-endian; the Motorola 68000, the IBM mainframes and the
PowerPC of classic Macintosh computers are big-endian (when Apple moved to Intel processors in 2006, it had to deal with the byte order of existing files and programs); ARM, PowerPC and MIPS are
bi-endian. The Internet protocols transmit numbers in big-endian order (the \emph{network byte order}). Many file formats therefore announce their byte order: a TIFF image starts with
\texttt{II} (Intel) or \texttt{MM} (Motorola), and a UTF-16 text file may begin with a byte-order mark.

\paragraph{Endianness of qubits.} Exactly the same ambiguity arises for qubits. A basis state of $n$ qubits is a string of bits, which can also be read as an integer $k$, the index of the corresponding
component of the state vector. The question is which qubit is the most significant bit of $k$.
\begin{itemize}
  \item \emph{Big-endian} (most textbooks, including Nielsen and Chuang \cite{NielsenChuang}, and these notes):
  \begin{equation}
    \ket{q_0q_1\dots q_{n-1}}=\ket{q_0}\otimes\ket{q_1}\otimes\dots\otimes\ket{q_{n-1}},\qquad k=\sum_{i=0}^{n-1}q_i2^{n-1-i}.
  \end{equation}
  \item \emph{Little-endian} (Qiskit):
  \begin{equation}
    \ket{q_{n-1}\dots q_1q_0}=\ket{q_{n-1}}\otimes\dots\otimes\ket{q_1}\otimes\ket{q_0},\qquad k=\sum_{i=0}^{n-1}q_i2^i .
  \end{equation}
  Qubit $q_i$ then has the weight $2^i$, as bit $i$ of an ordinary binary integer. This makes arithmetic circuits natural (and is consistent with the operation \texttt{1 << i} used in the
  Deutsch--Jozsa example), but it reverses the order of the tensor product.
\end{itemize}
\emph{Example 1: a single flipped qubit.} Apply an $X$ gate to qubit $q_0$ of a three-qubit register initially in $\ket{000}$:
\begin{lstlisting}
from qiskit import QuantumCircuit
from qiskit.quantum_info import Statevector

qc = QuantumCircuit(3)
qc.x(0)
print(Statevector(qc).probabilities_dict())   # {'001': 1.0}
\end{lstlisting}
Qiskit reports the outcome \texttt{'001'}, with $q_0$ as the rightmost bit, and the non-zero component of the state vector has index $k=1$. In the big-endian convention, the same state is written
$\ket{100}$ and has index $k=4$.

\emph{Example 2: the CNOT matrix.} For a CNOT with control $q_0$ and target $q_1$:
\begin{equation}
  \underbrace{\begin{pmatrix}1&0&0&0\\0&1&0&0\\0&0&0&1\\0&0&1&0\end{pmatrix}}_{\text{big-endian (these notes)}}\qquad
  \underbrace{\begin{pmatrix}1&0&0&0\\0&0&0&1\\0&0&1&0\\0&1&0&0\end{pmatrix}}_{\text{little-endian (Qiskit)}}
\end{equation}
In the Qiskit matrix, the control $q_0$ is the least significant bit, so the gate exchanges the basis states with indices $1=\ket{01}$ and $3=\ket{11}$ instead of 2 and 3. Both matrices describe
the same physical gate; only the labelling of the basis has changed. Formally, the two matrices are related by the SWAP permutation: $M_{\rm Qiskit}=\mathrm{SWAP}\,M\,\mathrm{SWAP}$.

\emph{Converting.} To compare Qiskit results with calculations done by hand, reverse the order of the qubits, with \texttt{QuantumCircuit.reverse\_bits()} or the method \texttt{reverse\_qargs()} of
the \texttt{Statevector} and \texttt{Operator} classes. For example, \texttt{Operator(qc.reverse\_bits())} gives the CNOT matrix in the big-endian form above. As in Lilliput, there is no right end
at which to break the egg. What matters is to know which convention a given book, article or program uses, and to convert carefully.

\section{Noisy simulation and fake backends}
To test a program under realistic conditions without access to a real quantum computer, Qiskit provides \emph{fake backends}: simulators that reproduce the qubit connectivity, the basis gates and the
noise of specific IBM devices, for example \texttt{FakeManilaV2} or \texttt{FakeMarrakesh} from \texttt{qiskit\_ibm\_runtime.fake\_provider}. The noise model contains, for each qubit and gate,
depolarizing errors, thermal relaxation with the measured $T_1$, $T_2$ (\cref{ch:open}), and readout errors.

\begin{keypoint}[: a back-of-the-envelope error estimate]
If a transpiled circuit has $N_1$ single-qubit gates with error $\varepsilon_1$, $N_2$ two-qubit gates with error $\varepsilon_2$ and $n$ measurements with error $\varepsilon_m$, the probability
that no error occurs is roughly
\begin{equation}
  P_{\rm success}\approx(1-\varepsilon_1)^{N_1}(1-\varepsilon_2)^{N_2}(1-\varepsilon_m)^{n}\approx\e^{-N_1\varepsilon_1-N_2\varepsilon_2-n\varepsilon_m}.
\end{equation}
With $\varepsilon_2\sim5\times10^{-3}$, a circuit becomes useless after a few hundred two-qubit gates. This is the first estimate to make when answering ``can it work in our case?''.
\end{keypoint}

\section{Higher-level tools}
Besides individual gates, Qiskit offers many higher-level building blocks:
\begin{itemize}
  \item \texttt{qiskit.circuit.library}: a library of ready-made circuits. It contains standard and generalized gates, arithmetic operations, Boolean logic, data-encoding circuits for quantum
  machine learning (e.g.\ the Pauli expansion \texttt{PauliFeatureMap} for supervised learning), state preparation (\texttt{Initialize}), phase and bit-flip oracles, the Grover operator,
  the QFT, and Hamiltonian evolution gates (\texttt{HamiltonianGate}, \texttt{PauliEvolutionGate}, which implements Trotter formulas \cref{eq:trotter}).
  \item \texttt{qiskit.quantum\_info}: utilities for quantum information. It covers quantum states and operators (e.g.\ the \texttt{Statevector} and \texttt{DensityMatrix} classes), partial traces,
  entropies, distances between states, fidelity (\cref{app:fidelity}), quantum channels (Kraus representation), and purity.
\end{itemize}

\section{The half adder}\label{sec:halfadder}
\subsection{The classical half adder}
The simplest arithmetic operation is the addition of two bits $x$ and $y$. The result can be $0$, $1$ or $2=10_2$, so it needs two output bits: the sum $S$ (weight 1) and the carry $C$ (weight 2):
\begin{center}
\begin{tabular}{@{}cc|cc@{}}\toprule
$x$&$y$&$C$&$S$\\ \midrule 0&0&0&0\\0&1&0&1\\1&0&0&1\\1&1&1&0\\ \bottomrule
\end{tabular}
\end{center}
The sum is the exclusive OR of the inputs and the carry is their AND:
\begin{equation}
  S=x\oplus y,\qquad C=x\wedge y.
\end{equation}
The circuit that implements these two functions is called a \emph{half adder}. It is only ``half'' of what is needed to add multi-bit numbers. When adding two binary numbers column by column,
each column must add three bits: the two digits and the carry coming from the previous column. A circuit that adds three bits is called a \emph{full adder}. It can be built from two half adders:
the first adds $x$ and $y$, the second adds the incoming carry to their sum, and the two partial carries are combined with an OR gate.

Adders are among the most important circuits of a classical computer. The arithmetic logic unit (ALU) of every processor is built around them. Through two's complement arithmetic, the same
adder also performs subtraction. Multiplication reduces to repeated shifts and additions, and adders are used constantly to compute memory addresses and increment the program counter.

\subsection{A quantum half adder}
A quantum circuit must be reversible. The classical half adder is not: the outputs $(C,S)=(0,1)$ do not tell us whether the input was $(0,1)$ or $(1,0)$. We therefore keep the inputs and write the
results into two additional qubits $s_0$ (sum) and $s_1$ (carry), both initialized to $\ket0$:
\begin{equation}
  \ket x\ket y\ket0\ket0\longrightarrow\ket x\ket y\ket{x\oplus y}\ket{x\wedge y}.
\end{equation}
\begin{itemize}
  \item The XOR is obtained with two CNOT gates targeting $s_0$: the first adds $x$, the second adds $y$ modulo 2, giving $s_0=x\oplus y$.
  \item The AND is obtained with a Toffoli gate (\texttt{ccx}), which flips the target $s_1$ only if both controls are $1$, giving $s_1=x\wedge y$.
\end{itemize}
\begin{lstlisting}
from qiskit import QuantumRegister, QuantumCircuit, transpile
from qiskit_aer import Aer

x = QuantumRegister(1, 'x')
y = QuantumRegister(1, 'y')
s = QuantumRegister(2, 's')   # s[0]: sum, s[1]: carry
qc = QuantumCircuit(x, y, s)

qc.h(x)                       # superposition of inputs 0 and 1
qc.h(y)
qc.barrier()
qc.cx(x, s[0])                # s0 = x
qc.cx(y, s[0])                # s0 = x XOR y
qc.ccx(x, y, s[1])            # s1 = x AND y
qc.barrier()
qc.measure_all()

simulator = Aer.get_backend('qasm_simulator')
counts = simulator.run(transpile(qc, simulator), shots=1000).result().get_counts()
\end{lstlisting}
\begin{figure}[h]\centering
\begin{quantikz}
\lstick{$x$} & \gate{H} & \slice{} & \ctrl{2} & \qw & \ctrl{3} & \slice{} & \meter{}\\
\lstick{$y$} & \gate{H} & \qw & \qw & \ctrl{1} & \ctrl{2} & \qw & \meter{}\\
\lstick{$s_0$} & \qw & \qw & \targ{} & \targ{} & \qw & \qw & \meter{}\\
\lstick{$s_1$} & \qw & \qw & \qw & \qw & \targ{} & \qw & \meter{}
\end{quantikz}
\caption{Quantum half adder. The Hadamard gates prepare the inputs in a superposition of all four values. Two CNOT gates compute the sum $s_0=x\oplus y$, and a Toffoli gate computes the carry $s_1=x\wedge y$.}
\label{fig:halfadder}
\end{figure}
Instead of a single input value, the Hadamard gates prepare the superposition of all four inputs. The circuit then evaluates the four additions simultaneously:
\begin{equation}
  \frac12\left(\ket{00}+\ket{01}+\ket{10}+\ket{11}\right)_{xy}\ket{00}_s\longrightarrow\frac12\sum_{x,y=0}^1\ket x\ket y\ket{x\oplus y}\ket{x\wedge y}.
\end{equation}
A run on the ideal simulator gives, for 1000 shots, approximately \texttt{\{'0000': 250, '0110': 250, '0101': 250, '1011': 250\}}. In Qiskit's little-endian convention the bit strings read
$s_1s_0yx$. For example, \texttt{'1011'} means $x=1$, $y=1$, sum $s_0=0$ and carry $s_1=1$, i.e.\ $1+1=10_2$. All four lines of the truth table appear, each with probability $1/4$. This illustrates
both the power and the limitation of quantum parallelism: the four additions are performed by a single run, but a measurement reveals only one of them, chosen at random. Useful quantum algorithms
must use interference to concentrate the probability on the desired answer before measuring.

\paragraph{Running on (simulated) hardware.} With a fake backend reproducing a real IBM processor, for example \texttt{FakeMarrakesh}, a few percent of the shots give outcomes that are not in the truth
table, such as \texttt{'1111'} or \texttt{'0111'}. The reason becomes clear after transpilation. The device has only one native two-qubit gate and limited connectivity, so the Toffoli gate alone is
decomposed into six two-qubit gates plus single-qubit rotations (\cref{sec:toffoli}), and SWAPs are added to bring the qubits together. The three logical gates of \cref{fig:halfadder} become about
15 two-qubit gates and nearly 60 single-qubit operations. With $\varepsilon_2\approx3\times10^{-3}$ and $\varepsilon_1\approx3\times10^{-4}$, the estimate of the previous section gives
$1-P_{\rm success}\approx15\times0.003+60\times0.0003\approx6\%$, in agreement with the observed error rate.

\subsection{Towards multi-bit addition: the ripple-carry adder}
To add two $n$-bit numbers $a$ and $b$, $n$ adders are chained as in addition by hand. The carry produced by the adder of bit $i$ is passed as an input to the adder of bit $i+1$, so the carry
``ripples'' from the least to the most significant bit. This \emph{ripple-carry adder} handles the incoming carries and adds numbers of arbitrary size, which a half adder cannot do. Its cost grows
only linearly with $n$.

Quantum versions of the ripple-carry adder were among the first quantum arithmetic circuits to be designed \cite{VBE1996}. They are better than a naive chain of half adders in an important respect:
the number of qubits. A naive reversible construction must keep every intermediate result, which consumes many ancilla qubits. The circuit of Cuccaro et al.\ \cite{Cuccaro2004} computes the sum in
place, $\ket a\ket b\to\ket a\ket{a+b}$, using a single ancilla qubit, by uncomputing the intermediate carries once they are no longer needed. Since qubits are the scarcest resource of present-day
quantum computers, this economy is essential. Both circuits are available in Qiskit as \texttt{VBERippleCarryAdder} and \texttt{CDKMRippleCarryAdder}, together with the single-bit
\texttt{HalfAdderGate} and \texttt{FullAdderGate}.

In classical processors, the ripple-carry adder is the simplest but also the slowest multi-bit adder, because the last bit must wait for the carry to propagate through all the previous ones. Faster
designs, such as carry-lookahead adders, compute the carries in parallel. The same idea leads to quantum adders of logarithmic depth, and to adders based on the quantum Fourier transform
(\texttt{DraperQFTAdder}): in the Fourier basis, adding $a$ is a set of single-qubit phase rotations (\cref{ch:shor}). Adders are the building blocks of modular exponentiation, the most expensive part
of Shor's algorithm.
